\subsection{Common-word realization on the three independent pair types}

\paragraph{D23.P8 (the three-pair commutator image is a direct product).}
For the matrices of D23.P6, the simultaneous evaluation map
\[
 \widehat F'\longrightarrow
 P_{12}'\times G_{34}'\times G_{45}'
\tag{D23.P8.1}
\]
is surjective. Here the factors are the derived groups of the
three corresponding pair images. In particular each of the $216$
reduced assignments in D23.P7 has a common commutator-word
preimage for these three pair values.

This statement does not yet prescribe the reversed mixed-pair
value, and it does not assert that the common word satisfies
the profinite gentle identities. These qualifications matter:
the reductions in D23.P7 used those identities for actual gentle
elements, whereas surjectivity in (D23.P8.1) concerns unrestricted
commutator words.

\begin{proof}
Write $G_{34}'=H_{34}'\times S_{34}$ and
$G_{45}'=H_{45}'\times S_{45}$ in the exact splitting bases of
D23.C9, where $|H_{34}'|=|H_{45}'|=27$ and
$S_{34}\simeq S_{45}\simeq\operatorname{SU}_3(3)$.

First, the two solvable full pair blocks $H_{34},H_{45}$ have
abelianization $C_4$ via determinant. With $-i$ as generator,
their two marked determinant phase pairs are $(2,1)$ for $x$
and $(1,1)$ for $y$. In any common abelian quotient of the two
blocks, if $a,b$ denote the images of their respective cyclic
abelianization generators, the common marking would impose
$2a=b$ and $a=b$. Hence $a=b=0$. If their simultaneous image
were a proper subdirect product, Goursat's lemma would give a
nontrivial common solvable quotient, whose abelianization would
be nontrivial. The displayed calculation excludes that.
Thus the full simultaneous image is $H_{34}\times H_{45}$,
and its derived image is $H_{34}'\times H_{45}'$.

Next the two simple factors have different kernels even on
$\widehat F'$. The following ordinary commutator word is a
small exact witness:
\[
 w=y^4x^2yx^{-1}y^{-1}x^{-1}y^{-4}xy^2x^2yxy^{-1}x^{-3}y^{-2}x^{-1}.
\tag{D23.P8.2}
\]
Its exponent sums are both zero. In the unitary block matrices
of D23.C9 it has values
\[
 w_{34}=I,\qquad
 w_{45}=\begin{pmatrix}5&7&4\\8&8&6\\5&7&5\end{pmatrix}\ne I,
\tag{D23.P8.3}
\]
using the same encoding $a+3b$ for $a+bi$ as the certificates.
The script \path{agents/pentagon/krammer_commonword_witness.py}
verifies the two values by literal word multiplication.
Since each factor is nonabelian simple, a subdirect product
of the two is either the whole product or the graph of an
isomorphism. The latter would give equal kernels, contrary
to (D23.P8.3). Thus the perfect joint image is
$S_{34}\times S_{45}$.

The standard simplicity of $\operatorname{SU}_3(3)$ is the
primary-source input used here. Its centre is trivial since
$\gcd(3,4)=1$; the GAP Character Table Library's
\emph{Simple Groups by Group Order}, row $U_3(3)$, lists its
order $6048$, and the corresponding group record identifies
it with $\operatorname{PSU}(3,3)$. This is the same cited
input used in D23.F14, not an inference from the order alone.

Finally $P_{12}'$ is solvable, and its abelianization has
order $3888/486=8$, by D23.P6. The group
$H_{34}'\times H_{45}'$ is a three-group. Any nontrivial
common quotient would be solvable and would have a nontrivial
abelian quotient, simultaneously a quotient of a two-group
and a three-group. This is impossible. Hence Goursat gives
the full product of these three solvable factors. It has no
nontrivial common quotient with the perfect product
$S_{34}\times S_{45}$, because a finite group both solvable
and perfect is trivial. A final application of Goursat proves
(D23.P8.1).

The image of the ordinary free commutator subgroup in any
finite target equals the image of its profinite closure.
Therefore ordinary commutator-word preimages can be chosen.
This proves the stated finite-word realization and its scope.
\end{proof}

\paragraph{Primary-source locators.}
\url{https://www.math.rwth-aachen.de/homes/Thomas.Breuer/ctbllib/ctbltoc/views/simplebyorder.html},
row $U_3(3)$;
\url{https://www.math.rwth-aachen.de/homes/sam/ctbllib/ctbltoc/data/U3(3).html},
``Group constructions in GAP'' entry $\operatorname{PSU}(3,3)$.
